
\documentclass{article}
\usepackage{colortbl}
\usepackage{makecell}
\usepackage{multirow}
\usepackage{supertabular}

\begin{document}

\newcounter{utterance}

\centering \large Interaction Transcript for game `dond', experiment `coop\_en', episode 7 with Qwen3{-}32B{-}t0.0.
\vspace{24pt}

{ \footnotesize  \setcounter{utterance}{1}
\setlength{\tabcolsep}{0pt}
\begin{supertabular}{c@{$\;$}|p{.15\linewidth}@{}p{.15\linewidth}p{.15\linewidth}p{.15\linewidth}p{.15\linewidth}p{.15\linewidth}}
    \# & $\;$A & \multicolumn{4}{c}{Game Master} & $\;\:$B\\
    \hline

    \theutterance \stepcounter{utterance}  
    & & \multicolumn{4}{p{0.6\linewidth}}{
        \cellcolor[rgb]{0.9,0.9,0.9}{
            \makecell[{{p{\linewidth}}}]{
                \texttt{\tiny{[P1$\langle$GM]}}
                \texttt{You are playing a cooperative negotiation game in which you have to agree on how to divide a set of items among you and another player.} \\
\\ 
\texttt{Rules:} \\
\texttt{(a) You and the other player are given a set of items. Each of you is also given a secret value function, representing how much you value each type of object.} \\
\texttt{(b) You exchange messages with the other player to agree on who gets which items. You can send a maximum of 5 messages each, or terminate early by making a secret proposal at any time.} \\
\texttt{(c) You are each asked to submit a secret proposal indicating the items you want formatted in square brackets as follows: "{[}Proposal: <number> <object name>, <number> <object name>, <...>{]}"} \\
\texttt{(d) If your proposals are complementary, i.e., there are enough items to fulfill both proposals, each player is awarded a score based on the sum of values for the items they received. Otherwise, both of you get zero points.} \\
\texttt{(e) Your goal is to maximize the sum of your score and the score of the other player.} \\
\\ 
\texttt{Let us start.} \\
\\ 
\texttt{The set of available items is:} \\
\\ 
\texttt{2 hooks, 2 plates, 1 bowl} \\
\\ 
\texttt{Your secret value function is:} \\
\\ 
\texttt{hook: 3, plate: 1, bowl: 2} \\
\\ 
\texttt{IMPORTANT: Your messages, unless it is the secret proposal, are directly transmitted to the other player, so do not include any response to the rules or text announcing your message. To make a secret proposal, use the indicated format. Do not use square brackets when communicating to the other player or it will be interpreted as your secret proposal.} \\
            }
        }
    }
    & \\ \\

    \theutterance \stepcounter{utterance}  
    & \multicolumn{4}{p{0.6\linewidth}}{
        \cellcolor[rgb]{0.9,0.9,0.9}{
            \makecell[{{p{\linewidth}}}]{
                \texttt{\tiny{[P1$\rangle$GM]}}
                \texttt{<think>} \\
\texttt{Okay, let's see. I need to divide the items between me and another player. The items are 2 hooks, 2 plates, and 1 bowl. My value function is hook:3, plate:1, bowl:2. The other player's value function is secret, so I don't know what they value. The goal is to maximize the total score for both of us. } \\
\\ 
\texttt{First, I should figure out what the other player might value. Since I don't have their value function, I need to communicate to find out. But I can only send 5 messages each. Let me think about possible strategies.} \\
\\ 
\texttt{Maybe I should start by asking them what they value most. But since the messages are direct, I can't ask questions like that. Wait, the rules say that messages are exchanged, but I can't include any response to the rules or text announcing my message. So I need to be careful with my messages. } \\
\\ 
\texttt{Alternatively, maybe I can propose a division and see if they accept. But since the proposals are secret, if I make a secret proposal, they can make their own. If the proposals are complementary, we get points. Otherwise, zero. So the key is to make sure our proposals don't overlap.} \\
\\ 
\texttt{Let me think about possible allocations. Let's consider the total value for each item. For me, hooks are high value (3 each), plates low (1), bowl medium (2). The other player's values are unknown. But to maximize the total sum, we need to allocate items in a way that both of us get the items we value most.} \\
\\ 
\texttt{But without knowing their values, I need to make assumptions. Maybe the other player values the bowl more? Or maybe they value plates more? Let's consider possible scenarios.} \\
\\ 
\texttt{If the other player values the bowl highly, then maybe I should take the hooks and plates, and let them take the bowl. But I have to consider their possible value. Alternatively, if they value hooks, maybe we can split the hooks. But there are only two hooks. If I take one hook, they take one, but then we can each get 3 points from that. But I also have to consider other items.} \\
\\ 
\texttt{Wait, the total items are 2 hooks, 2 plates, 1 bowl. So total of 5 items. If I take two hooks, that's 6 points for me, but leaves plates and a bowl. If the other player values} \\
            }
        }
    }
    & & \\ \\

    \theutterance \stepcounter{utterance}  
    & & & \multicolumn{2}{p{0.3\linewidth}}{
        \cellcolor[rgb]{0.9,0.9,0.9}{
            \makecell[{{p{\linewidth}}}]{
                \texttt{\tiny{[GM$|$GM]}}
                \texttt{continue} \\
            }
        }
    }
    & & \\ \\

    \theutterance \stepcounter{utterance}  
    & & & \multicolumn{2}{p{0.3\linewidth}}{
        \cellcolor[rgb]{0.9,0.9,0.9}{
            \makecell[{{p{\linewidth}}}]{
                \texttt{\tiny{[GM$|$GM]}}
                \texttt{<think>} \\
\texttt{Okay, let's see. I need to divide the items between me and another player. The items are 2 hooks, 2 plates, and 1 bowl. My value function is hook:3, plate:1, bowl:2. The other player's value function is secret, so I don't know what they value. The goal is to maximize the total score for both of us. } \\
\\ 
\texttt{First, I should figure out what the other player might value. Since I don't have their value function, I need to communicate to find out. But I can only send 5 messages each. Let me think about possible strategies.} \\
\\ 
\texttt{Maybe I should start by asking them what they value most. But since the messages are direct, I can't ask questions like that. Wait, the rules say that messages are exchanged, but I can't include any response to the rules or text announcing my message. So I need to be careful with my messages. } \\
\\ 
\texttt{Alternatively, maybe I can propose a division and see if they accept. But since the proposals are secret, if I make a secret proposal, they can make their own. If the proposals are complementary, we get points. Otherwise, zero. So the key is to make sure our proposals don't overlap.} \\
\\ 
\texttt{Let me think about possible allocations. Let's consider the total value for each item. For me, hooks are high value (3 each), plates low (1), bowl medium (2). The other player's values are unknown. But to maximize the total sum, we need to allocate items in a way that both of us get the items we value most.} \\
\\ 
\texttt{But without knowing their values, I need to make assumptions. Maybe the other player values the bowl more? Or maybe they value plates more? Let's consider possible scenarios.} \\
\\ 
\texttt{If the other player values the bowl highly, then maybe I should take the hooks and plates, and let them take the bowl. But I have to consider their possible value. Alternatively, if they value hooks, maybe we can split the hooks. But there are only two hooks. If I take one hook, they take one, but then we can each get 3 points from that. But I also have to consider other items.} \\
\\ 
\texttt{Wait, the total items are 2 hooks, 2 plates, 1 bowl. So total of 5 items. If I take two hooks, that's 6 points for me, but leaves plates and a bowl. If the other player values} \\
            }
        }
    }
    & & \\ \\

    \theutterance \stepcounter{utterance}  
    & & \multicolumn{4}{p{0.6\linewidth}}{
        \cellcolor[rgb]{0.9,0.9,0.9}{
            \makecell[{{p{\linewidth}}}]{
                \texttt{\tiny{[GM$\rangle$P2]}}
                \texttt{You are playing a cooperative negotiation game in which you have to agree on how to divide a set of items among you and another player.} \\
\\ 
\texttt{Rules:} \\
\texttt{(a) You and the other player are given a set of items. Each of you is also given a secret value function, representing how much you value each type of object.} \\
\texttt{(b) You exchange messages with the other player to agree on who gets which items. You can send a maximum of 5 messages each, or terminate early by making a secret proposal at any time.} \\
\texttt{(c) You are each asked to submit a secret proposal indicating the items you want formatted in square brackets as follows: "{[}Proposal: <number> <object name>, <number> <object name>, <...>{]}"} \\
\texttt{(d) If your proposals are complementary, i.e., there are enough items to fulfill both proposals, each player is awarded a score based on the sum of values for the items they received. Otherwise, both of you get zero points.} \\
\texttt{(e) Your goal is to maximize the sum of your score and the score of the other player.} \\
\\ 
\texttt{Let us start.} \\
\\ 
\texttt{The set of available items is:} \\
\\ 
\texttt{2 hooks, 2 plates, 1 bowl} \\
\\ 
\texttt{Your secret value function is:} \\
\\ 
\texttt{hook: 1, plate: 0, bowl: 8} \\
\\ 
\texttt{IMPORTANT: Your messages, unless it is the secret proposal, are directly transmitted to the other player, so do not include any response to the rules or text announcing your message. To make a secret proposal, use the indicated format. Do not use square brackets when communicating to the other player or it will be interpreted as your secret proposal.} \\
\\ 
\\ 
\texttt{<think>} \\
\texttt{Okay, let's see. I need to divide the items between me and another player. The items are 2 hooks, 2 plates, and 1 bowl. My value function is hook:3, plate:1, bowl:2. The other player's value function is secret, so I don't know what they value. The goal is to maximize the total score for both of us. } \\
\\ 
\texttt{First, I should figure out what the other player might value. Since I don't have their value function, I need to communicate to find out. But I can only send 5 messages each. Let me think about possible strategies.} \\
\\ 
\texttt{Maybe I should start by asking them what they value most. But since the messages are direct, I can't ask questions like that. Wait, the rules say that messages are exchanged, but I can't include any response to the rules or text announcing my message. So I need to be careful with my messages. } \\
\\ 
\texttt{Alternatively, maybe I can propose a division and see if they accept. But since the proposals are secret, if I make a secret proposal, they can make their own. If the proposals are complementary, we get points. Otherwise, zero. So the key is to make sure our proposals don't overlap.} \\
\\ 
\texttt{Let me think about possible allocations. Let's consider the total value for each item. For me, hooks are high value (3 each), plates low (1), bowl medium (2). The other player's values are unknown. But to maximize the total sum, we need to allocate items in a way that both of us get the items we value most.} \\
\\ 
\texttt{But without knowing their values, I need to make assumptions. Maybe the other player values the bowl more? Or maybe they value plates more? Let's consider possible scenarios.} \\
\\ 
\texttt{If the other player values the bowl highly, then maybe I should take the hooks and plates, and let them take the bowl. But I have to consider their possible value. Alternatively, if they value hooks, maybe we can split the hooks. But there are only two hooks. If I take one hook, they take one, but then we can each get 3 points from that. But I also have to consider other items.} \\
\\ 
\texttt{Wait, the total items are 2 hooks, 2 plates, 1 bowl. So total of 5 items. If I take two hooks, that's 6 points for me, but leaves plates and a bowl. If the other player values} \\
            }
        }
    }
    & \\ \\

    \theutterance \stepcounter{utterance}  
    & & & \multicolumn{4}{p{0.6\linewidth}}{
        \cellcolor[rgb]{0.9,0.9,0.9}{
            \makecell[{{p{\linewidth}}}]{
                \texttt{\tiny{[GM$\langle$P2]}}
                \texttt{<think>} \\
\texttt{Okay, let's tackle this problem step by step. The goal is to divide the items (2 hooks, 2 plates, 1 bowl) between me and another player in a way that maximizes the total score for both of us. My value function is hook:3, plate:1, bowl:2. The other player's values are secret, so I need to figure out a strategy that works regardless of their preferences.} \\
\\ 
\texttt{First, I need to consider possible allocations. Since the other player's values are unknown, I should aim for a division that is likely to be complementary to their potential proposals. Let's list out all possible ways to split the items and see which ones might maximize the total.} \\
\\ 
\texttt{Option 1: I take both hooks and the bowl. That gives me 3+3+2=8 points. The remaining items are 2 plates. If the other player values plates highly, they might take both, but since plates are worth 1 each to me, maybe they value them more. However, if they take both plates, their total would depend on their values. But if they need the bowl, they might not get it if I take it. This could lead to conflict.} \\
\\ 
\texttt{Option 2: I take both hooks and one plate. That's 3+3+1=7 points. The remaining items are 1 plate and 1 bowl. If the other player wants the bowl, they can take it along with the remaining plate. Their total would depend on their values. If they value the bowl highly, this might work. But if they also want hooks, which I already took, there's no conflict here since I took both.} \\
\\ 
\texttt{Option 3: Split the hooks. I take one hook and the bowl, giving me 3+2=5. The other player takes the other hook and maybe the plates. Their total would depend on their values. If they value hooks more, they get 3 plus whatever they get from plates. But I might not get the maximum here.} \\
\\ 
\texttt{Option 4: I take the bowl and one plate. That's 2+1=3. The other player could take the hooks and the remaining plate. If they value hooks highly, this would be good for them. My total is low, but maybe their total is higher, leading to a better overall sum.} \\
\\ 
\texttt{But since my own values are higher for hooks, I should prioritize getting hooks. However, if the other player} \\
            }
        }
    }
     \\ \\

    \theutterance \stepcounter{utterance}  
    & & & \multicolumn{2}{p{0.3\linewidth}}{
        \cellcolor[rgb]{0.9,0.9,0.9}{
            \makecell[{{p{\linewidth}}}]{
                \texttt{\tiny{[GM$|$GM]}}
                \texttt{continue} \\
            }
        }
    }
    & & \\ \\

    \theutterance \stepcounter{utterance}  
    & & & \multicolumn{2}{p{0.3\linewidth}}{
        \cellcolor[rgb]{0.9,0.9,0.9}{
            \makecell[{{p{\linewidth}}}]{
                \texttt{\tiny{[GM$|$GM]}}
                \texttt{<think>} \\
\texttt{Okay, let's tackle this problem step by step. The goal is to divide the items (2 hooks, 2 plates, 1 bowl) between me and another player in a way that maximizes the total score for both of us. My value function is hook:3, plate:1, bowl:2. The other player's values are secret, so I need to figure out a strategy that works regardless of their preferences.} \\
\\ 
\texttt{First, I need to consider possible allocations. Since the other player's values are unknown, I should aim for a division that is likely to be complementary to their potential proposals. Let's list out all possible ways to split the items and see which ones might maximize the total.} \\
\\ 
\texttt{Option 1: I take both hooks and the bowl. That gives me 3+3+2=8 points. The remaining items are 2 plates. If the other player values plates highly, they might take both, but since plates are worth 1 each to me, maybe they value them more. However, if they take both plates, their total would depend on their values. But if they need the bowl, they might not get it if I take it. This could lead to conflict.} \\
\\ 
\texttt{Option 2: I take both hooks and one plate. That's 3+3+1=7 points. The remaining items are 1 plate and 1 bowl. If the other player wants the bowl, they can take it along with the remaining plate. Their total would depend on their values. If they value the bowl highly, this might work. But if they also want hooks, which I already took, there's no conflict here since I took both.} \\
\\ 
\texttt{Option 3: Split the hooks. I take one hook and the bowl, giving me 3+2=5. The other player takes the other hook and maybe the plates. Their total would depend on their values. If they value hooks more, they get 3 plus whatever they get from plates. But I might not get the maximum here.} \\
\\ 
\texttt{Option 4: I take the bowl and one plate. That's 2+1=3. The other player could take the hooks and the remaining plate. If they value hooks highly, this would be good for them. My total is low, but maybe their total is higher, leading to a better overall sum.} \\
\\ 
\texttt{But since my own values are higher for hooks, I should prioritize getting hooks. However, if the other player} \\
            }
        }
    }
    & & \\ \\

    \theutterance \stepcounter{utterance}  
    & & \multicolumn{4}{p{0.6\linewidth}}{
        \cellcolor[rgb]{0.9,0.9,0.9}{
            \makecell[{{p{\linewidth}}}]{
                \texttt{\tiny{[P1$\langle$GM]}}
                \texttt{<think>} \\
\texttt{Okay, let's tackle this problem step by step. The goal is to divide the items (2 hooks, 2 plates, 1 bowl) between me and another player in a way that maximizes the total score for both of us. My value function is hook:3, plate:1, bowl:2. The other player's values are secret, so I need to figure out a strategy that works regardless of their preferences.} \\
\\ 
\texttt{First, I need to consider possible allocations. Since the other player's values are unknown, I should aim for a division that is likely to be complementary to their potential proposals. Let's list out all possible ways to split the items and see which ones might maximize the total.} \\
\\ 
\texttt{Option 1: I take both hooks and the bowl. That gives me 3+3+2=8 points. The remaining items are 2 plates. If the other player values plates highly, they might take both, but since plates are worth 1 each to me, maybe they value them more. However, if they take both plates, their total would depend on their values. But if they need the bowl, they might not get it if I take it. This could lead to conflict.} \\
\\ 
\texttt{Option 2: I take both hooks and one plate. That's 3+3+1=7 points. The remaining items are 1 plate and 1 bowl. If the other player wants the bowl, they can take it along with the remaining plate. Their total would depend on their values. If they value the bowl highly, this might work. But if they also want hooks, which I already took, there's no conflict here since I took both.} \\
\\ 
\texttt{Option 3: Split the hooks. I take one hook and the bowl, giving me 3+2=5. The other player takes the other hook and maybe the plates. Their total would depend on their values. If they value hooks more, they get 3 plus whatever they get from plates. But I might not get the maximum here.} \\
\\ 
\texttt{Option 4: I take the bowl and one plate. That's 2+1=3. The other player could take the hooks and the remaining plate. If they value hooks highly, this would be good for them. My total is low, but maybe their total is higher, leading to a better overall sum.} \\
\\ 
\texttt{But since my own values are higher for hooks, I should prioritize getting hooks. However, if the other player} \\
            }
        }
    }
    & \\ \\

    \theutterance \stepcounter{utterance}  
    & \multicolumn{4}{p{0.6\linewidth}}{
        \cellcolor[rgb]{0.9,0.9,0.9}{
            \makecell[{{p{\linewidth}}}]{
                \texttt{\tiny{[P1$\rangle$GM]}}
                \texttt{<think>} \\
\texttt{Okay, let's see. I need to figure out how to divide the items between me and the other player to maximize the total score. The items are 2 hooks, 2 plates, and 1 bowl. My values are hook:3, plate:1, bowl:2. The other player's values are secret, so I don't know what they value. The key is to make sure our proposals are complementary so we both get points.} \\
\\ 
\texttt{First, I should consider possible allocations. Since I value hooks the most, maybe I should take as many hooks as possible. But if the other player also values hooks, we might both want them, leading to a conflict. So maybe I should take one hook and let the other player take the other. But there are two hooks, so splitting them could work. Let's think about different scenarios.} \\
\\ 
\texttt{If I take both hooks, that's 6 points for me. Then the remaining items are 2 plates and 1 bowl. The other player could take the bowl and one plate, or both plates. If they value the bowl more, they might take it. But if they value plates more, they might take both. However, if they take both plates, that's 2 points for them (if they value plates at 1 each), but maybe they value them higher. But since I don't know their values, it's risky.} \\
\\ 
\texttt{Alternatively, if I take one hook and the bowl, that's 3 + 2 = 5 points. Then the other player can take the other hook and the two plates. If they value hooks highly, they get 3 + (plates' value). If they value plates more, they might get more from the plates. But again, without knowing their values, it's hard to say.} \\
\\ 
\texttt{Another option is to take one hook and one plate. That's 3 + 1 = 4 points. Then the other player can take the other hook, the remaining plate, and the bowl. If they value the bowl and hooks, that could be good for them. But again, depends on their values.} \\
\\ 
\texttt{Wait, the goal is to maximize the total score for both. So maybe the best approach is to allocate items in a way that both of us get the items we value most. Since I don't know their values, I need to make an assumption. Let's assume the other player might value the bowl more than the plates. If I take the} \\
            }
        }
    }
    & & \\ \\

    \theutterance \stepcounter{utterance}  
    & & & \multicolumn{2}{p{0.3\linewidth}}{
        \cellcolor[rgb]{0.9,0.9,0.9}{
            \makecell[{{p{\linewidth}}}]{
                \texttt{\tiny{[GM$|$GM]}}
                \texttt{continue} \\
            }
        }
    }
    & & \\ \\

    \theutterance \stepcounter{utterance}  
    & & & \multicolumn{2}{p{0.3\linewidth}}{
        \cellcolor[rgb]{0.9,0.9,0.9}{
            \makecell[{{p{\linewidth}}}]{
                \texttt{\tiny{[GM$|$GM]}}
                \texttt{<think>} \\
\texttt{Okay, let's see. I need to figure out how to divide the items between me and the other player to maximize the total score. The items are 2 hooks, 2 plates, and 1 bowl. My values are hook:3, plate:1, bowl:2. The other player's values are secret, so I don't know what they value. The key is to make sure our proposals are complementary so we both get points.} \\
\\ 
\texttt{First, I should consider possible allocations. Since I value hooks the most, maybe I should take as many hooks as possible. But if the other player also values hooks, we might both want them, leading to a conflict. So maybe I should take one hook and let the other player take the other. But there are two hooks, so splitting them could work. Let's think about different scenarios.} \\
\\ 
\texttt{If I take both hooks, that's 6 points for me. Then the remaining items are 2 plates and 1 bowl. The other player could take the bowl and one plate, or both plates. If they value the bowl more, they might take it. But if they value plates more, they might take both. However, if they take both plates, that's 2 points for them (if they value plates at 1 each), but maybe they value them higher. But since I don't know their values, it's risky.} \\
\\ 
\texttt{Alternatively, if I take one hook and the bowl, that's 3 + 2 = 5 points. Then the other player can take the other hook and the two plates. If they value hooks highly, they get 3 + (plates' value). If they value plates more, they might get more from the plates. But again, without knowing their values, it's hard to say.} \\
\\ 
\texttt{Another option is to take one hook and one plate. That's 3 + 1 = 4 points. Then the other player can take the other hook, the remaining plate, and the bowl. If they value the bowl and hooks, that could be good for them. But again, depends on their values.} \\
\\ 
\texttt{Wait, the goal is to maximize the total score for both. So maybe the best approach is to allocate items in a way that both of us get the items we value most. Since I don't know their values, I need to make an assumption. Let's assume the other player might value the bowl more than the plates. If I take the} \\
            }
        }
    }
    & & \\ \\

    \theutterance \stepcounter{utterance}  
    & & \multicolumn{4}{p{0.6\linewidth}}{
        \cellcolor[rgb]{0.9,0.9,0.9}{
            \makecell[{{p{\linewidth}}}]{
                \texttt{\tiny{[GM$\rangle$P2]}}
                \texttt{<think>} \\
\texttt{Okay, let's see. I need to figure out how to divide the items between me and the other player to maximize the total score. The items are 2 hooks, 2 plates, and 1 bowl. My values are hook:3, plate:1, bowl:2. The other player's values are secret, so I don't know what they value. The key is to make sure our proposals are complementary so we both get points.} \\
\\ 
\texttt{First, I should consider possible allocations. Since I value hooks the most, maybe I should take as many hooks as possible. But if the other player also values hooks, we might both want them, leading to a conflict. So maybe I should take one hook and let the other player take the other. But there are two hooks, so splitting them could work. Let's think about different scenarios.} \\
\\ 
\texttt{If I take both hooks, that's 6 points for me. Then the remaining items are 2 plates and 1 bowl. The other player could take the bowl and one plate, or both plates. If they value the bowl more, they might take it. But if they value plates more, they might take both. However, if they take both plates, that's 2 points for them (if they value plates at 1 each), but maybe they value them higher. But since I don't know their values, it's risky.} \\
\\ 
\texttt{Alternatively, if I take one hook and the bowl, that's 3 + 2 = 5 points. Then the other player can take the other hook and the two plates. If they value hooks highly, they get 3 + (plates' value). If they value plates more, they might get more from the plates. But again, without knowing their values, it's hard to say.} \\
\\ 
\texttt{Another option is to take one hook and one plate. That's 3 + 1 = 4 points. Then the other player can take the other hook, the remaining plate, and the bowl. If they value the bowl and hooks, that could be good for them. But again, depends on their values.} \\
\\ 
\texttt{Wait, the goal is to maximize the total score for both. So maybe the best approach is to allocate items in a way that both of us get the items we value most. Since I don't know their values, I need to make an assumption. Let's assume the other player might value the bowl more than the plates. If I take the} \\
            }
        }
    }
    & \\ \\

    \theutterance \stepcounter{utterance}  
    & & & \multicolumn{4}{p{0.6\linewidth}}{
        \cellcolor[rgb]{0.9,0.9,0.9}{
            \makecell[{{p{\linewidth}}}]{
                \texttt{\tiny{[GM$\langle$P2]}}
                \texttt{<think>} \\
\texttt{Okay, let's try to figure this out. The user is trying to divide the items between themselves and another player, with the goal of maximizing the total score for both. The items are 2 hooks, 2 plates, and 1 bowl. The user's values are hook:3, plate:1, bowl:2. The other player's values are secret.} \\
\\ 
\texttt{First, the user needs to consider possible allocations that would allow both parties to get items they value without overlapping. Since the other player's values are unknown, the user has to make educated guesses. The key is to find a division where the user's proposal and the other player's proposal don't conflict.} \\
\\ 
\texttt{The user's highest value items are hooks, followed by the bowl, then plates. So ideally, the user would want as many hooks as possible. But if the other player also values hooks, there might be a conflict. However, since there are two hooks, maybe splitting them could work. Let's explore possible allocations.} \\
\\ 
\texttt{Option 1: User takes both hooks (3+3=6 points) and maybe the bowl (2 points). That would leave the plates for the other player. If the other player values plates more, they could take both. But if the other player values the bowl more, they might want it, leading to a conflict.} \\
\\ 
\texttt{Option 2: User takes one hook (3) and the bowl (2), totaling 5. The other player could take the remaining hook (3) and the two plates (1+1=2). If the other player values hooks and plates, this could work. The total would be 5 (user) + 3+2=5 (other) = 10.} \\
\\ 
\texttt{Option 3: User takes one hook (3) and one plate (1), totaling 4. The other player takes the other hook (3), the bowl (2), and the remaining plate (1). Total would be 4 (user) + 3+2+1=6 (other) = 10. But the user's total is lower here.} \\
\\ 
\texttt{Option 4: User takes the bowl (2) and both plates (1+1=2), totaling 4. The other player takes both hooks (3+3=6). Total would be 4 + 6 = 10. But again, the user's score is lower.} \\
\\ 
\texttt{The maximum total in these scenarios is 10.} \\
            }
        }
    }
     \\ \\

    \theutterance \stepcounter{utterance}  
    & & & \multicolumn{2}{p{0.3\linewidth}}{
        \cellcolor[rgb]{0.9,0.9,0.9}{
            \makecell[{{p{\linewidth}}}]{
                \texttt{\tiny{[GM$|$GM]}}
                \texttt{continue} \\
            }
        }
    }
    & & \\ \\

    \theutterance \stepcounter{utterance}  
    & & & \multicolumn{2}{p{0.3\linewidth}}{
        \cellcolor[rgb]{0.9,0.9,0.9}{
            \makecell[{{p{\linewidth}}}]{
                \texttt{\tiny{[GM$|$GM]}}
                \texttt{<think>} \\
\texttt{Okay, let's try to figure this out. The user is trying to divide the items between themselves and another player, with the goal of maximizing the total score for both. The items are 2 hooks, 2 plates, and 1 bowl. The user's values are hook:3, plate:1, bowl:2. The other player's values are secret.} \\
\\ 
\texttt{First, the user needs to consider possible allocations that would allow both parties to get items they value without overlapping. Since the other player's values are unknown, the user has to make educated guesses. The key is to find a division where the user's proposal and the other player's proposal don't conflict.} \\
\\ 
\texttt{The user's highest value items are hooks, followed by the bowl, then plates. So ideally, the user would want as many hooks as possible. But if the other player also values hooks, there might be a conflict. However, since there are two hooks, maybe splitting them could work. Let's explore possible allocations.} \\
\\ 
\texttt{Option 1: User takes both hooks (3+3=6 points) and maybe the bowl (2 points). That would leave the plates for the other player. If the other player values plates more, they could take both. But if the other player values the bowl more, they might want it, leading to a conflict.} \\
\\ 
\texttt{Option 2: User takes one hook (3) and the bowl (2), totaling 5. The other player could take the remaining hook (3) and the two plates (1+1=2). If the other player values hooks and plates, this could work. The total would be 5 (user) + 3+2=5 (other) = 10.} \\
\\ 
\texttt{Option 3: User takes one hook (3) and one plate (1), totaling 4. The other player takes the other hook (3), the bowl (2), and the remaining plate (1). Total would be 4 (user) + 3+2+1=6 (other) = 10. But the user's total is lower here.} \\
\\ 
\texttt{Option 4: User takes the bowl (2) and both plates (1+1=2), totaling 4. The other player takes both hooks (3+3=6). Total would be 4 + 6 = 10. But again, the user's score is lower.} \\
\\ 
\texttt{The maximum total in these scenarios is 10.} \\
            }
        }
    }
    & & \\ \\

    \theutterance \stepcounter{utterance}  
    & & \multicolumn{4}{p{0.6\linewidth}}{
        \cellcolor[rgb]{0.9,0.9,0.9}{
            \makecell[{{p{\linewidth}}}]{
                \texttt{\tiny{[P1$\langle$GM]}}
                \texttt{<think>} \\
\texttt{Okay, let's try to figure this out. The user is trying to divide the items between themselves and another player, with the goal of maximizing the total score for both. The items are 2 hooks, 2 plates, and 1 bowl. The user's values are hook:3, plate:1, bowl:2. The other player's values are secret.} \\
\\ 
\texttt{First, the user needs to consider possible allocations that would allow both parties to get items they value without overlapping. Since the other player's values are unknown, the user has to make educated guesses. The key is to find a division where the user's proposal and the other player's proposal don't conflict.} \\
\\ 
\texttt{The user's highest value items are hooks, followed by the bowl, then plates. So ideally, the user would want as many hooks as possible. But if the other player also values hooks, there might be a conflict. However, since there are two hooks, maybe splitting them could work. Let's explore possible allocations.} \\
\\ 
\texttt{Option 1: User takes both hooks (3+3=6 points) and maybe the bowl (2 points). That would leave the plates for the other player. If the other player values plates more, they could take both. But if the other player values the bowl more, they might want it, leading to a conflict.} \\
\\ 
\texttt{Option 2: User takes one hook (3) and the bowl (2), totaling 5. The other player could take the remaining hook (3) and the two plates (1+1=2). If the other player values hooks and plates, this could work. The total would be 5 (user) + 3+2=5 (other) = 10.} \\
\\ 
\texttt{Option 3: User takes one hook (3) and one plate (1), totaling 4. The other player takes the other hook (3), the bowl (2), and the remaining plate (1). Total would be 4 (user) + 3+2+1=6 (other) = 10. But the user's total is lower here.} \\
\\ 
\texttt{Option 4: User takes the bowl (2) and both plates (1+1=2), totaling 4. The other player takes both hooks (3+3=6). Total would be 4 + 6 = 10. But again, the user's score is lower.} \\
\\ 
\texttt{The maximum total in these scenarios is 10.} \\
            }
        }
    }
    & \\ \\

    \theutterance \stepcounter{utterance}  
    & \multicolumn{4}{p{0.6\linewidth}}{
        \cellcolor[rgb]{0.9,0.9,0.9}{
            \makecell[{{p{\linewidth}}}]{
                \texttt{\tiny{[P1$\rangle$GM]}}
                \texttt{<think>} \\
\texttt{Okay, let's try to figure out the best way to divide the items. The user's values are hook:3, plate:1, bowl:2. The other player's values are unknown, but we need to maximize the total score for both. The items are 2 hooks, 2 plates, 1 bowl.} \\
\\ 
\texttt{First, I need to consider possible allocations where both players can get items they value without overlapping. Since the other player's values are secret, I have to make assumptions. Let's think about different scenarios.} \\
\\ 
\texttt{If I take both hooks (3+3=6) and the bowl (2), that's 8 points for me. The remaining items are 2 plates. If the other player values plates more, they could take both, but since plates are only worth 1 each to me, maybe they value them higher. However, if they need the bowl, they might not get it if I take it. This could lead to a conflict.} \\
\\ 
\texttt{Alternatively, if I take one hook (3) and the bowl (2), that's 5 points. The other player can take the remaining hook (3) and the two plates (1+1=2). If they value hooks and plates, their total would be 5. The combined total is 10. That seems good.} \\
\\ 
\texttt{Another option: I take one hook (3) and one plate (1), totaling 4. The other player takes the other hook (3), the bowl (2), and the remaining plate (1). Their total would be 6. Combined total is 10 again. But my score is lower here.} \\
\\ 
\texttt{If I take the bowl (2) and both plates (1+1=2), that's 4 points. The other player takes both hooks (3+3=6). Total combined is 10. But again, my score is lower.} \\
\\ 
\texttt{The maximum combined total in these scenarios is 10. To achieve this, the user needs to take either one hook and the bowl, or one hook and one plate, allowing the other player to take the remaining items. However, since the user's own value is higher for hooks, taking one hook and the bowl gives them 5 points, which is better than 4. The other player would then take the remaining hook and plates, which might be better for them if they value hooks and plates.} \\
\\ 
\texttt{But wait, there's another possibility. If the other player} \\
            }
        }
    }
    & & \\ \\

    \theutterance \stepcounter{utterance}  
    & & & \multicolumn{2}{p{0.3\linewidth}}{
        \cellcolor[rgb]{0.9,0.9,0.9}{
            \makecell[{{p{\linewidth}}}]{
                \texttt{\tiny{[GM$|$GM]}}
                \texttt{continue} \\
            }
        }
    }
    & & \\ \\

    \theutterance \stepcounter{utterance}  
    & & & \multicolumn{2}{p{0.3\linewidth}}{
        \cellcolor[rgb]{0.9,0.9,0.9}{
            \makecell[{{p{\linewidth}}}]{
                \texttt{\tiny{[GM$|$GM]}}
                \texttt{<think>} \\
\texttt{Okay, let's try to figure out the best way to divide the items. The user's values are hook:3, plate:1, bowl:2. The other player's values are unknown, but we need to maximize the total score for both. The items are 2 hooks, 2 plates, 1 bowl.} \\
\\ 
\texttt{First, I need to consider possible allocations where both players can get items they value without overlapping. Since the other player's values are secret, I have to make assumptions. Let's think about different scenarios.} \\
\\ 
\texttt{If I take both hooks (3+3=6) and the bowl (2), that's 8 points for me. The remaining items are 2 plates. If the other player values plates more, they could take both, but since plates are only worth 1 each to me, maybe they value them higher. However, if they need the bowl, they might not get it if I take it. This could lead to a conflict.} \\
\\ 
\texttt{Alternatively, if I take one hook (3) and the bowl (2), that's 5 points. The other player can take the remaining hook (3) and the two plates (1+1=2). If they value hooks and plates, their total would be 5. The combined total is 10. That seems good.} \\
\\ 
\texttt{Another option: I take one hook (3) and one plate (1), totaling 4. The other player takes the other hook (3), the bowl (2), and the remaining plate (1). Their total would be 6. Combined total is 10 again. But my score is lower here.} \\
\\ 
\texttt{If I take the bowl (2) and both plates (1+1=2), that's 4 points. The other player takes both hooks (3+3=6). Total combined is 10. But again, my score is lower.} \\
\\ 
\texttt{The maximum combined total in these scenarios is 10. To achieve this, the user needs to take either one hook and the bowl, or one hook and one plate, allowing the other player to take the remaining items. However, since the user's own value is higher for hooks, taking one hook and the bowl gives them 5 points, which is better than 4. The other player would then take the remaining hook and plates, which might be better for them if they value hooks and plates.} \\
\\ 
\texttt{But wait, there's another possibility. If the other player} \\
            }
        }
    }
    & & \\ \\

    \theutterance \stepcounter{utterance}  
    & & \multicolumn{4}{p{0.6\linewidth}}{
        \cellcolor[rgb]{0.9,0.9,0.9}{
            \makecell[{{p{\linewidth}}}]{
                \texttt{\tiny{[GM$\rangle$P2]}}
                \texttt{<think>} \\
\texttt{Okay, let's try to figure out the best way to divide the items. The user's values are hook:3, plate:1, bowl:2. The other player's values are unknown, but we need to maximize the total score for both. The items are 2 hooks, 2 plates, 1 bowl.} \\
\\ 
\texttt{First, I need to consider possible allocations where both players can get items they value without overlapping. Since the other player's values are secret, I have to make assumptions. Let's think about different scenarios.} \\
\\ 
\texttt{If I take both hooks (3+3=6) and the bowl (2), that's 8 points for me. The remaining items are 2 plates. If the other player values plates more, they could take both, but since plates are only worth 1 each to me, maybe they value them higher. However, if they need the bowl, they might not get it if I take it. This could lead to a conflict.} \\
\\ 
\texttt{Alternatively, if I take one hook (3) and the bowl (2), that's 5 points. The other player can take the remaining hook (3) and the two plates (1+1=2). If they value hooks and plates, their total would be 5. The combined total is 10. That seems good.} \\
\\ 
\texttt{Another option: I take one hook (3) and one plate (1), totaling 4. The other player takes the other hook (3), the bowl (2), and the remaining plate (1). Their total would be 6. Combined total is 10 again. But my score is lower here.} \\
\\ 
\texttt{If I take the bowl (2) and both plates (1+1=2), that's 4 points. The other player takes both hooks (3+3=6). Total combined is 10. But again, my score is lower.} \\
\\ 
\texttt{The maximum combined total in these scenarios is 10. To achieve this, the user needs to take either one hook and the bowl, or one hook and one plate, allowing the other player to take the remaining items. However, since the user's own value is higher for hooks, taking one hook and the bowl gives them 5 points, which is better than 4. The other player would then take the remaining hook and plates, which might be better for them if they value hooks and plates.} \\
\\ 
\texttt{But wait, there's another possibility. If the other player} \\
            }
        }
    }
    & \\ \\

    \theutterance \stepcounter{utterance}  
    & & & \multicolumn{4}{p{0.6\linewidth}}{
        \cellcolor[rgb]{0.9,0.9,0.9}{
            \makecell[{{p{\linewidth}}}]{
                \texttt{\tiny{[GM$\langle$P2]}}
                \texttt{<think>} \\
\texttt{Okay, let's try to figure out the best way to divide the items between me and the other player. The items are 2 hooks, 2 plates, and 1 bowl. My values are hook:3, plate:1, bowl:2. The other player's values are secret, so I need to make an allocation that maximizes the total score for both of us.} \\
\\ 
\texttt{First, I need to consider possible allocations where our proposals don't overlap. Since I don't know their values, I have to make some assumptions. Let's think about different scenarios.} \\
\\ 
\texttt{If I take both hooks (3+3=6) and the bowl (2), that's 8 points for me. The remaining items are 2 plates. If the other player values plates more, they could take both, but since plates are only worth 1 each to me, maybe they value them higher. However, if they need the bowl, they might not get it if I take it. This could lead to a conflict.} \\
\\ 
\texttt{Alternatively, if I take one hook (3) and the bowl (2), that's 5 points. The other player can take the remaining hook (3) and the two plates (1+1=2). If they value hooks and plates, their total would be 5. The combined total is 10. That seems good.} \\
\\ 
\texttt{Another option: I take one hook (3) and one plate (1), totaling 4. The other player takes the other hook (3), the bowl (2), and the remaining plate (1). Their total would be 6. Combined total is 10 again. But my score is lower here.} \\
\\ 
\texttt{If I take the bowl (2) and both plates (1+1=2), that's 4 points. The other player takes both hooks (3+3=6). Total combined is 10. But again, my score is lower.} \\
\\ 
\texttt{The maximum combined total in these scenarios is 10. To achieve this, the user needs to take either one hook and the bowl, or one hook and one plate, allowing the other player to take the remaining items. However, since the user's own value is higher for hooks, taking one hook and the bowl gives them 5 points, which is better than 4. The other player would then take the remaining hook and plates, which might be better for them if they value hooks and plates.} \\
\\ 
\texttt{But wait, there} \\
            }
        }
    }
     \\ \\

    \theutterance \stepcounter{utterance}  
    & & & \multicolumn{2}{p{0.3\linewidth}}{
        \cellcolor[rgb]{0.9,0.9,0.9}{
            \makecell[{{p{\linewidth}}}]{
                \texttt{\tiny{[GM$|$GM]}}
                \texttt{continue} \\
            }
        }
    }
    & & \\ \\

    \theutterance \stepcounter{utterance}  
    & & & \multicolumn{2}{p{0.3\linewidth}}{
        \cellcolor[rgb]{0.9,0.9,0.9}{
            \makecell[{{p{\linewidth}}}]{
                \texttt{\tiny{[GM$|$GM]}}
                \texttt{<think>} \\
\texttt{Okay, let's try to figure out the best way to divide the items between me and the other player. The items are 2 hooks, 2 plates, and 1 bowl. My values are hook:3, plate:1, bowl:2. The other player's values are secret, so I need to make an allocation that maximizes the total score for both of us.} \\
\\ 
\texttt{First, I need to consider possible allocations where our proposals don't overlap. Since I don't know their values, I have to make some assumptions. Let's think about different scenarios.} \\
\\ 
\texttt{If I take both hooks (3+3=6) and the bowl (2), that's 8 points for me. The remaining items are 2 plates. If the other player values plates more, they could take both, but since plates are only worth 1 each to me, maybe they value them higher. However, if they need the bowl, they might not get it if I take it. This could lead to a conflict.} \\
\\ 
\texttt{Alternatively, if I take one hook (3) and the bowl (2), that's 5 points. The other player can take the remaining hook (3) and the two plates (1+1=2). If they value hooks and plates, their total would be 5. The combined total is 10. That seems good.} \\
\\ 
\texttt{Another option: I take one hook (3) and one plate (1), totaling 4. The other player takes the other hook (3), the bowl (2), and the remaining plate (1). Their total would be 6. Combined total is 10 again. But my score is lower here.} \\
\\ 
\texttt{If I take the bowl (2) and both plates (1+1=2), that's 4 points. The other player takes both hooks (3+3=6). Total combined is 10. But again, my score is lower.} \\
\\ 
\texttt{The maximum combined total in these scenarios is 10. To achieve this, the user needs to take either one hook and the bowl, or one hook and one plate, allowing the other player to take the remaining items. However, since the user's own value is higher for hooks, taking one hook and the bowl gives them 5 points, which is better than 4. The other player would then take the remaining hook and plates, which might be better for them if they value hooks and plates.} \\
\\ 
\texttt{But wait, there} \\
            }
        }
    }
    & & \\ \\

    \theutterance \stepcounter{utterance}  
    & & \multicolumn{4}{p{0.6\linewidth}}{
        \cellcolor[rgb]{0.9,0.9,0.9}{
            \makecell[{{p{\linewidth}}}]{
                \texttt{\tiny{[P1$\langle$GM]}}
                \texttt{<think>} \\
\texttt{Okay, let's try to figure out the best way to divide the items between me and the other player. The items are 2 hooks, 2 plates, and 1 bowl. My values are hook:3, plate:1, bowl:2. The other player's values are secret, so I need to make an allocation that maximizes the total score for both of us.} \\
\\ 
\texttt{First, I need to consider possible allocations where our proposals don't overlap. Since I don't know their values, I have to make some assumptions. Let's think about different scenarios.} \\
\\ 
\texttt{If I take both hooks (3+3=6) and the bowl (2), that's 8 points for me. The remaining items are 2 plates. If the other player values plates more, they could take both, but since plates are only worth 1 each to me, maybe they value them higher. However, if they need the bowl, they might not get it if I take it. This could lead to a conflict.} \\
\\ 
\texttt{Alternatively, if I take one hook (3) and the bowl (2), that's 5 points. The other player can take the remaining hook (3) and the two plates (1+1=2). If they value hooks and plates, their total would be 5. The combined total is 10. That seems good.} \\
\\ 
\texttt{Another option: I take one hook (3) and one plate (1), totaling 4. The other player takes the other hook (3), the bowl (2), and the remaining plate (1). Their total would be 6. Combined total is 10 again. But my score is lower here.} \\
\\ 
\texttt{If I take the bowl (2) and both plates (1+1=2), that's 4 points. The other player takes both hooks (3+3=6). Total combined is 10. But again, my score is lower.} \\
\\ 
\texttt{The maximum combined total in these scenarios is 10. To achieve this, the user needs to take either one hook and the bowl, or one hook and one plate, allowing the other player to take the remaining items. However, since the user's own value is higher for hooks, taking one hook and the bowl gives them 5 points, which is better than 4. The other player would then take the remaining hook and plates, which might be better for them if they value hooks and plates.} \\
\\ 
\texttt{But wait, there} \\
            }
        }
    }
    & \\ \\

    \theutterance \stepcounter{utterance}  
    & \multicolumn{4}{p{0.6\linewidth}}{
        \cellcolor[rgb]{0.9,0.9,0.9}{
            \makecell[{{p{\linewidth}}}]{
                \texttt{\tiny{[P1$\rangle$GM]}}
                \texttt{<think>} \\
\texttt{Okay, let's try to figure out the best way to divide the items between me and the other player. The items are 2 hooks, 2 plates, and 1 bowl. My values are hook:3, plate:1, bowl:2. The other player's values are secret, so I need to make an allocation that maximizes the total score for both of us.} \\
\\ 
\texttt{First, I need to consider possible allocations where our proposals don't overlap. Since I don't know their values, I have to make some assumptions. Let's think about different scenarios.} \\
\\ 
\texttt{If I take both hooks (3+3=6) and the bowl (2), that's 8 points for me. The remaining items are 2 plates. If the other player values plates more, they could take both, but since plates are only worth 1 each to me, maybe they value them higher. However, if they need the bowl, they might not get it if I take it. This could lead to a conflict.} \\
\\ 
\texttt{Alternatively, if I take one hook (3) and the bowl (2), that's 5 points. The other player can take the remaining hook (3) and the two plates (1+1=2). If they value hooks and plates, their total would be 5. The combined total is 10. That seems good.} \\
\\ 
\texttt{Another option: I take one hook (3) and one plate (1), totaling 4. The other player takes the other hook (3), the bowl (2), and the remaining plate (1). Their total would be 6. Combined total is 10 again. But my score is lower here.} \\
\\ 
\texttt{If I take the bowl (2) and both plates (1+1=2), that's 4 points. The other player takes both hooks (3+3=6). Total combined is 10. But again, my score is lower.} \\
\\ 
\texttt{The maximum combined total in these scenarios is 10. To achieve this, the user needs to take either one hook and the bowl, or one hook and one plate, allowing the other player to take the remaining items. However, since the user's own value is higher for hooks, taking one hook and the bowl gives them 5 points, which is better than 4. The other player would then take the remaining hook and plates, which might be better for them if they value hooks and plates.} \\
\\ 
\texttt{But wait, there} \\
            }
        }
    }
    & & \\ \\

    \theutterance \stepcounter{utterance}  
    & & & \multicolumn{2}{p{0.3\linewidth}}{
        \cellcolor[rgb]{0.9,0.9,0.9}{
            \makecell[{{p{\linewidth}}}]{
                \texttt{\tiny{[GM$|$GM]}}
                \texttt{continue} \\
            }
        }
    }
    & & \\ \\

    \theutterance \stepcounter{utterance}  
    & & & \multicolumn{2}{p{0.3\linewidth}}{
        \cellcolor[rgb]{0.9,0.9,0.9}{
            \makecell[{{p{\linewidth}}}]{
                \texttt{\tiny{[GM$|$GM]}}
                \texttt{<think>} \\
\texttt{Okay, let's try to figure out the best way to divide the items between me and the other player. The items are 2 hooks, 2 plates, and 1 bowl. My values are hook:3, plate:1, bowl:2. The other player's values are secret, so I need to make an allocation that maximizes the total score for both of us.} \\
\\ 
\texttt{First, I need to consider possible allocations where our proposals don't overlap. Since I don't know their values, I have to make some assumptions. Let's think about different scenarios.} \\
\\ 
\texttt{If I take both hooks (3+3=6) and the bowl (2), that's 8 points for me. The remaining items are 2 plates. If the other player values plates more, they could take both, but since plates are only worth 1 each to me, maybe they value them higher. However, if they need the bowl, they might not get it if I take it. This could lead to a conflict.} \\
\\ 
\texttt{Alternatively, if I take one hook (3) and the bowl (2), that's 5 points. The other player can take the remaining hook (3) and the two plates (1+1=2). If they value hooks and plates, their total would be 5. The combined total is 10. That seems good.} \\
\\ 
\texttt{Another option: I take one hook (3) and one plate (1), totaling 4. The other player takes the other hook (3), the bowl (2), and the remaining plate (1). Their total would be 6. Combined total is 10 again. But my score is lower here.} \\
\\ 
\texttt{If I take the bowl (2) and both plates (1+1=2), that's 4 points. The other player takes both hooks (3+3=6). Total combined is 10. But again, my score is lower.} \\
\\ 
\texttt{The maximum combined total in these scenarios is 10. To achieve this, the user needs to take either one hook and the bowl, or one hook and one plate, allowing the other player to take the remaining items. However, since the user's own value is higher for hooks, taking one hook and the bowl gives them 5 points, which is better than 4. The other player would then take the remaining hook and plates, which might be better for them if they value hooks and plates.} \\
\\ 
\texttt{But wait, there} \\
            }
        }
    }
    & & \\ \\

    \theutterance \stepcounter{utterance}  
    & & \multicolumn{4}{p{0.6\linewidth}}{
        \cellcolor[rgb]{0.9,0.9,0.9}{
            \makecell[{{p{\linewidth}}}]{
                \texttt{\tiny{[GM$\rangle$P2]}}
                \texttt{<think>} \\
\texttt{Okay, let's try to figure out the best way to divide the items between me and the other player. The items are 2 hooks, 2 plates, and 1 bowl. My values are hook:3, plate:1, bowl:2. The other player's values are secret, so I need to make an allocation that maximizes the total score for both of us.} \\
\\ 
\texttt{First, I need to consider possible allocations where our proposals don't overlap. Since I don't know their values, I have to make some assumptions. Let's think about different scenarios.} \\
\\ 
\texttt{If I take both hooks (3+3=6) and the bowl (2), that's 8 points for me. The remaining items are 2 plates. If the other player values plates more, they could take both, but since plates are only worth 1 each to me, maybe they value them higher. However, if they need the bowl, they might not get it if I take it. This could lead to a conflict.} \\
\\ 
\texttt{Alternatively, if I take one hook (3) and the bowl (2), that's 5 points. The other player can take the remaining hook (3) and the two plates (1+1=2). If they value hooks and plates, their total would be 5. The combined total is 10. That seems good.} \\
\\ 
\texttt{Another option: I take one hook (3) and one plate (1), totaling 4. The other player takes the other hook (3), the bowl (2), and the remaining plate (1). Their total would be 6. Combined total is 10 again. But my score is lower here.} \\
\\ 
\texttt{If I take the bowl (2) and both plates (1+1=2), that's 4 points. The other player takes both hooks (3+3=6). Total combined is 10. But again, my score is lower.} \\
\\ 
\texttt{The maximum combined total in these scenarios is 10. To achieve this, the user needs to take either one hook and the bowl, or one hook and one plate, allowing the other player to take the remaining items. However, since the user's own value is higher for hooks, taking one hook and the bowl gives them 5 points, which is better than 4. The other player would then take the remaining hook and plates, which might be better for them if they value hooks and plates.} \\
\\ 
\texttt{But wait, there} \\
            }
        }
    }
    & \\ \\

    \theutterance \stepcounter{utterance}  
    & & & \multicolumn{4}{p{0.6\linewidth}}{
        \cellcolor[rgb]{0.9,0.9,0.9}{
            \makecell[{{p{\linewidth}}}]{
                \texttt{\tiny{[GM$\langle$P2]}}
                \texttt{<think>} \\
\texttt{Okay, let's try to figure out the best way to divide the items between me and the other player. The items are 2 hooks, 2 plates, and 1 bowl. My values are hook:3, plate:1, bowl:2. The other player's values are secret, so I need to make an allocation that maximizes the total score for both of us.} \\
\\ 
\texttt{First, I need to consider possible allocations where our proposals don't overlap. Since I don't know their values, I have to make some assumptions. Let's think about different scenarios.} \\
\\ 
\texttt{If I take both hooks (3+3=6) and the bowl (2), that's 8 points for me. The remaining items are 2 plates. If the other player values plates more, they could take both, but since plates are only worth 1 each to me, maybe they value them higher. However, if they need the bowl, they might not get it if I take it. This could lead to a conflict.} \\
\\ 
\texttt{Alternatively, if I take one hook (3) and the bowl (2), that's 5 points. The other player can take the remaining hook (3) and the two plates (1+1=2). If they value hooks and plates, their total would be 5. The combined total is 10. That seems good.} \\
\\ 
\texttt{Another option: I take one hook (3) and one plate (1), totaling 4. The other player takes the other hook (3), the bowl (2), and the remaining plate (1). Their total would be 6. Combined total is 10 again. But my score is lower here.} \\
\\ 
\texttt{If I take the bowl (2) and both plates (1+1=2), that's 4 points. The other player takes both hooks (3+3=6). Total combined is 10. But again, my score is lower.} \\
\\ 
\texttt{The maximum combined total in these scenarios is 10. To achieve this, the user needs to take either one hook and the bowl, or one hook and one plate, allowing the other player to take the remaining items. However, since the user's own value is higher for hooks, taking one hook and the bowl gives them 5 points, which is better than 4. The other player would then take the remaining hook and plates, which might be better for them if they value hooks and plates.} \\
\\ 
\texttt{But wait, there} \\
            }
        }
    }
     \\ \\

    \theutterance \stepcounter{utterance}  
    & & & \multicolumn{2}{p{0.3\linewidth}}{
        \cellcolor[rgb]{0.9,0.9,0.9}{
            \makecell[{{p{\linewidth}}}]{
                \texttt{\tiny{[GM$|$GM]}}
                \texttt{continue} \\
            }
        }
    }
    & & \\ \\

    \theutterance \stepcounter{utterance}  
    & & & \multicolumn{2}{p{0.3\linewidth}}{
        \cellcolor[rgb]{0.9,0.9,0.9}{
            \makecell[{{p{\linewidth}}}]{
                \texttt{\tiny{[GM$|$GM]}}
                \texttt{<think>} \\
\texttt{Okay, let's try to figure out the best way to divide the items between me and the other player. The items are 2 hooks, 2 plates, and 1 bowl. My values are hook:3, plate:1, bowl:2. The other player's values are secret, so I need to make an allocation that maximizes the total score for both of us.} \\
\\ 
\texttt{First, I need to consider possible allocations where our proposals don't overlap. Since I don't know their values, I have to make some assumptions. Let's think about different scenarios.} \\
\\ 
\texttt{If I take both hooks (3+3=6) and the bowl (2), that's 8 points for me. The remaining items are 2 plates. If the other player values plates more, they could take both, but since plates are only worth 1 each to me, maybe they value them higher. However, if they need the bowl, they might not get it if I take it. This could lead to a conflict.} \\
\\ 
\texttt{Alternatively, if I take one hook (3) and the bowl (2), that's 5 points. The other player can take the remaining hook (3) and the two plates (1+1=2). If they value hooks and plates, their total would be 5. The combined total is 10. That seems good.} \\
\\ 
\texttt{Another option: I take one hook (3) and one plate (1), totaling 4. The other player takes the other hook (3), the bowl (2), and the remaining plate (1). Their total would be 6. Combined total is 10 again. But my score is lower here.} \\
\\ 
\texttt{If I take the bowl (2) and both plates (1+1=2), that's 4 points. The other player takes both hooks (3+3=6). Total combined is 10. But again, my score is lower.} \\
\\ 
\texttt{The maximum combined total in these scenarios is 10. To achieve this, the user needs to take either one hook and the bowl, or one hook and one plate, allowing the other player to take the remaining items. However, since the user's own value is higher for hooks, taking one hook and the bowl gives them 5 points, which is better than 4. The other player would then take the remaining hook and plates, which might be better for them if they value hooks and plates.} \\
\\ 
\texttt{But wait, there} \\
            }
        }
    }
    & & \\ \\

    \theutterance \stepcounter{utterance}  
    & & \multicolumn{4}{p{0.6\linewidth}}{
        \cellcolor[rgb]{0.9,0.9,0.9}{
            \makecell[{{p{\linewidth}}}]{
                \texttt{\tiny{[P1$\langle$GM]}}
                \texttt{<think>} \\
\texttt{Okay, let's try to figure out the best way to divide the items between me and the other player. The items are 2 hooks, 2 plates, and 1 bowl. My values are hook:3, plate:1, bowl:2. The other player's values are secret, so I need to make an allocation that maximizes the total score for both of us.} \\
\\ 
\texttt{First, I need to consider possible allocations where our proposals don't overlap. Since I don't know their values, I have to make some assumptions. Let's think about different scenarios.} \\
\\ 
\texttt{If I take both hooks (3+3=6) and the bowl (2), that's 8 points for me. The remaining items are 2 plates. If the other player values plates more, they could take both, but since plates are only worth 1 each to me, maybe they value them higher. However, if they need the bowl, they might not get it if I take it. This could lead to a conflict.} \\
\\ 
\texttt{Alternatively, if I take one hook (3) and the bowl (2), that's 5 points. The other player can take the remaining hook (3) and the two plates (1+1=2). If they value hooks and plates, their total would be 5. The combined total is 10. That seems good.} \\
\\ 
\texttt{Another option: I take one hook (3) and one plate (1), totaling 4. The other player takes the other hook (3), the bowl (2), and the remaining plate (1). Their total would be 6. Combined total is 10 again. But my score is lower here.} \\
\\ 
\texttt{If I take the bowl (2) and both plates (1+1=2), that's 4 points. The other player takes both hooks (3+3=6). Total combined is 10. But again, my score is lower.} \\
\\ 
\texttt{The maximum combined total in these scenarios is 10. To achieve this, the user needs to take either one hook and the bowl, or one hook and one plate, allowing the other player to take the remaining items. However, since the user's own value is higher for hooks, taking one hook and the bowl gives them 5 points, which is better than 4. The other player would then take the remaining hook and plates, which might be better for them if they value hooks and plates.} \\
\\ 
\texttt{But wait, there} \\
            }
        }
    }
    & \\ \\

    \theutterance \stepcounter{utterance}  
    & \multicolumn{4}{p{0.6\linewidth}}{
        \cellcolor[rgb]{0.9,0.9,0.9}{
            \makecell[{{p{\linewidth}}}]{
                \texttt{\tiny{[P1$\rangle$GM]}}
                \texttt{<think>} \\
\texttt{Okay, let's try to figure out the best way to divide the items between me and the other player. The items are 2 hooks, 2 plates, and 1 bowl. My values are hook:3, plate:1, bowl:2. The other player's values are secret, so I need to make an allocation that maximizes the total score for both of us.} \\
\\ 
\texttt{First, I need to consider possible allocations where our proposals don't overlap. Since I don't know their values, I have to make some assumptions. Let's think about different scenarios.} \\
\\ 
\texttt{If I take both hooks (3+3=6) and the bowl (2), that's 8 points for me. The remaining items are 2 plates. If the other player values plates more, they could take both, but since plates are only worth 1 each to me, maybe they value them higher. However, if they need the bowl, they might not get it if I take it. This could lead to a conflict.} \\
\\ 
\texttt{Alternatively, if I take one hook (3) and the bowl (2), that's 5 points. The other player can take the remaining hook (3) and the two plates (1+1=2). If they value hooks and plates, their total would be 5. The combined total is 10. That seems good.} \\
\\ 
\texttt{Another option: I take one hook (3) and one plate (1), totaling 4. The other player takes the other hook (3), the bowl (2), and the remaining plate (1). Their total would be 6. Combined total is 10 again. But my score is lower here.} \\
\\ 
\texttt{If I take the bowl (2) and both plates (1+1=2), that's 4 points. The other player takes both hooks (3+3=6). Total combined is 10. But again, my score is lower.} \\
\\ 
\texttt{The maximum combined total in these scenarios is 10. To achieve this, the user needs to take either one hook and the bowl, or one hook and one plate, allowing the other player to take the remaining items. However, since the user's own value is higher for hooks, taking one hook and the bowl gives them 5 points, which is better than 4. The other player would then take the remaining hook and plates, which might be better for them if they value hooks and plates.} \\
\\ 
\texttt{But wait, there} \\
            }
        }
    }
    & & \\ \\

    \theutterance \stepcounter{utterance}  
    & & & \multicolumn{2}{p{0.3\linewidth}}{
        \cellcolor[rgb]{0.9,0.9,0.9}{
            \makecell[{{p{\linewidth}}}]{
                \texttt{\tiny{[GM$|$GM]}}
                \texttt{continue} \\
            }
        }
    }
    & & \\ \\

    \theutterance \stepcounter{utterance}  
    & & & \multicolumn{2}{p{0.3\linewidth}}{
        \cellcolor[rgb]{0.9,0.9,0.9}{
            \makecell[{{p{\linewidth}}}]{
                \texttt{\tiny{[GM$|$GM]}}
                \texttt{<think>} \\
\texttt{Okay, let's try to figure out the best way to divide the items between me and the other player. The items are 2 hooks, 2 plates, and 1 bowl. My values are hook:3, plate:1, bowl:2. The other player's values are secret, so I need to make an allocation that maximizes the total score for both of us.} \\
\\ 
\texttt{First, I need to consider possible allocations where our proposals don't overlap. Since I don't know their values, I have to make some assumptions. Let's think about different scenarios.} \\
\\ 
\texttt{If I take both hooks (3+3=6) and the bowl (2), that's 8 points for me. The remaining items are 2 plates. If the other player values plates more, they could take both, but since plates are only worth 1 each to me, maybe they value them higher. However, if they need the bowl, they might not get it if I take it. This could lead to a conflict.} \\
\\ 
\texttt{Alternatively, if I take one hook (3) and the bowl (2), that's 5 points. The other player can take the remaining hook (3) and the two plates (1+1=2). If they value hooks and plates, their total would be 5. The combined total is 10. That seems good.} \\
\\ 
\texttt{Another option: I take one hook (3) and one plate (1), totaling 4. The other player takes the other hook (3), the bowl (2), and the remaining plate (1). Their total would be 6. Combined total is 10 again. But my score is lower here.} \\
\\ 
\texttt{If I take the bowl (2) and both plates (1+1=2), that's 4 points. The other player takes both hooks (3+3=6). Total combined is 10. But again, my score is lower.} \\
\\ 
\texttt{The maximum combined total in these scenarios is 10. To achieve this, the user needs to take either one hook and the bowl, or one hook and one plate, allowing the other player to take the remaining items. However, since the user's own value is higher for hooks, taking one hook and the bowl gives them 5 points, which is better than 4. The other player would then take the remaining hook and plates, which might be better for them if they value hooks and plates.} \\
\\ 
\texttt{But wait, there} \\
            }
        }
    }
    & & \\ \\

    \theutterance \stepcounter{utterance}  
    & & \multicolumn{4}{p{0.6\linewidth}}{
        \cellcolor[rgb]{0.9,0.9,0.9}{
            \makecell[{{p{\linewidth}}}]{
                \texttt{\tiny{[GM$\rangle$P2]}}
                \texttt{<think>} \\
\texttt{Okay, let's try to figure out the best way to divide the items between me and the other player. The items are 2 hooks, 2 plates, and 1 bowl. My values are hook:3, plate:1, bowl:2. The other player's values are secret, so I need to make an allocation that maximizes the total score for both of us.} \\
\\ 
\texttt{First, I need to consider possible allocations where our proposals don't overlap. Since I don't know their values, I have to make some assumptions. Let's think about different scenarios.} \\
\\ 
\texttt{If I take both hooks (3+3=6) and the bowl (2), that's 8 points for me. The remaining items are 2 plates. If the other player values plates more, they could take both, but since plates are only worth 1 each to me, maybe they value them higher. However, if they need the bowl, they might not get it if I take it. This could lead to a conflict.} \\
\\ 
\texttt{Alternatively, if I take one hook (3) and the bowl (2), that's 5 points. The other player can take the remaining hook (3) and the two plates (1+1=2). If they value hooks and plates, their total would be 5. The combined total is 10. That seems good.} \\
\\ 
\texttt{Another option: I take one hook (3) and one plate (1), totaling 4. The other player takes the other hook (3), the bowl (2), and the remaining plate (1). Their total would be 6. Combined total is 10 again. But my score is lower here.} \\
\\ 
\texttt{If I take the bowl (2) and both plates (1+1=2), that's 4 points. The other player takes both hooks (3+3=6). Total combined is 10. But again, my score is lower.} \\
\\ 
\texttt{The maximum combined total in these scenarios is 10. To achieve this, the user needs to take either one hook and the bowl, or one hook and one plate, allowing the other player to take the remaining items. However, since the user's own value is higher for hooks, taking one hook and the bowl gives them 5 points, which is better than 4. The other player would then take the remaining hook and plates, which might be better for them if they value hooks and plates.} \\
\\ 
\texttt{But wait, there} \\
            }
        }
    }
    & \\ \\

    \theutterance \stepcounter{utterance}  
    & & & \multicolumn{4}{p{0.6\linewidth}}{
        \cellcolor[rgb]{0.9,0.9,0.9}{
            \makecell[{{p{\linewidth}}}]{
                \texttt{\tiny{[GM$\langle$P2]}}
                \texttt{<think>} \\
\texttt{Okay, let's try to figure out the best way to divide the items between me and the other player. The items are 2 hooks, 2 plates, and 1 bowl. My values are hook:3, plate:1, bowl:2. The other player's values are secret, so I need to make an allocation that maximizes the total score for both of us.} \\
\\ 
\texttt{First, I need to consider possible allocations where our proposals don't overlap. Since I don't know their values, I have to make some assumptions. Let's think about different scenarios.} \\
\\ 
\texttt{If I take both hooks (3+3=6) and the bowl (2), that's 8 points for me. The remaining items are 2 plates. If the other player values plates more, they could take both, but since plates are only worth 1 each to me, maybe they value them higher. However, if they need the bowl, they might not get it if I take it. This could lead to a conflict.} \\
\\ 
\texttt{Alternatively, if I take one hook (3) and the bowl (2), that's 5 points. The other player can take the remaining hook (3) and the two plates (1+1=2). If they value hooks and plates, their total would be 5. The combined total is 10. That seems good.} \\
\\ 
\texttt{Another option: I take one hook (3) and one plate (1), totaling 4. The other player takes the other hook (3), the bowl (2), and the remaining plate (1). Their total would be 6. Combined total is 10 again. But my score is lower here.} \\
\\ 
\texttt{If I take the bowl (2) and both plates (1+1=2), that's 4 points. The other player takes both hooks (3+3=6). Total combined is 10. But again, my score is lower.} \\
\\ 
\texttt{The maximum combined total in these scenarios is 10. To achieve this, the user needs to take either one hook and the bowl, or one hook and one plate, allowing the other player to take the remaining items. However, since the user's own value is higher for hooks, taking one hook and the bowl gives them 5 points, which is better than 4. The other player would then take the remaining hook and plates, which might be better for them if they value hooks and plates.} \\
\\ 
\texttt{But wait, there} \\
            }
        }
    }
     \\ \\

    \theutterance \stepcounter{utterance}  
    & & & \multicolumn{2}{p{0.3\linewidth}}{
        \cellcolor[rgb]{0.9,0.9,0.9}{
            \makecell[{{p{\linewidth}}}]{
                \texttt{\tiny{[GM$|$GM]}}
                \texttt{continue} \\
            }
        }
    }
    & & \\ \\

    \theutterance \stepcounter{utterance}  
    & & & \multicolumn{2}{p{0.3\linewidth}}{
        \cellcolor[rgb]{0.9,0.9,0.9}{
            \makecell[{{p{\linewidth}}}]{
                \texttt{\tiny{[GM$|$GM]}}
                \texttt{<think>} \\
\texttt{Okay, let's try to figure out the best way to divide the items between me and the other player. The items are 2 hooks, 2 plates, and 1 bowl. My values are hook:3, plate:1, bowl:2. The other player's values are secret, so I need to make an allocation that maximizes the total score for both of us.} \\
\\ 
\texttt{First, I need to consider possible allocations where our proposals don't overlap. Since I don't know their values, I have to make some assumptions. Let's think about different scenarios.} \\
\\ 
\texttt{If I take both hooks (3+3=6) and the bowl (2), that's 8 points for me. The remaining items are 2 plates. If the other player values plates more, they could take both, but since plates are only worth 1 each to me, maybe they value them higher. However, if they need the bowl, they might not get it if I take it. This could lead to a conflict.} \\
\\ 
\texttt{Alternatively, if I take one hook (3) and the bowl (2), that's 5 points. The other player can take the remaining hook (3) and the two plates (1+1=2). If they value hooks and plates, their total would be 5. The combined total is 10. That seems good.} \\
\\ 
\texttt{Another option: I take one hook (3) and one plate (1), totaling 4. The other player takes the other hook (3), the bowl (2), and the remaining plate (1). Their total would be 6. Combined total is 10 again. But my score is lower here.} \\
\\ 
\texttt{If I take the bowl (2) and both plates (1+1=2), that's 4 points. The other player takes both hooks (3+3=6). Total combined is 10. But again, my score is lower.} \\
\\ 
\texttt{The maximum combined total in these scenarios is 10. To achieve this, the user needs to take either one hook and the bowl, or one hook and one plate, allowing the other player to take the remaining items. However, since the user's own value is higher for hooks, taking one hook and the bowl gives them 5 points, which is better than 4. The other player would then take the remaining hook and plates, which might be better for them if they value hooks and plates.} \\
\\ 
\texttt{But wait, there} \\
            }
        }
    }
    & & \\ \\

    \theutterance \stepcounter{utterance}  
    & & \multicolumn{4}{p{0.6\linewidth}}{
        \cellcolor[rgb]{0.9,0.9,0.9}{
            \makecell[{{p{\linewidth}}}]{
                \texttt{\tiny{[P1$\langle$GM]}}
                \texttt{<think>} \\
\texttt{Okay, let's try to figure out the best way to divide the items between me and the other player. The items are 2 hooks, 2 plates, and 1 bowl. My values are hook:3, plate:1, bowl:2. The other player's values are secret, so I need to make an allocation that maximizes the total score for both of us.} \\
\\ 
\texttt{First, I need to consider possible allocations where our proposals don't overlap. Since I don't know their values, I have to make some assumptions. Let's think about different scenarios.} \\
\\ 
\texttt{If I take both hooks (3+3=6) and the bowl (2), that's 8 points for me. The remaining items are 2 plates. If the other player values plates more, they could take both, but since plates are only worth 1 each to me, maybe they value them higher. However, if they need the bowl, they might not get it if I take it. This could lead to a conflict.} \\
\\ 
\texttt{Alternatively, if I take one hook (3) and the bowl (2), that's 5 points. The other player can take the remaining hook (3) and the two plates (1+1=2). If they value hooks and plates, their total would be 5. The combined total is 10. That seems good.} \\
\\ 
\texttt{Another option: I take one hook (3) and one plate (1), totaling 4. The other player takes the other hook (3), the bowl (2), and the remaining plate (1). Their total would be 6. Combined total is 10 again. But my score is lower here.} \\
\\ 
\texttt{If I take the bowl (2) and both plates (1+1=2), that's 4 points. The other player takes both hooks (3+3=6). Total combined is 10. But again, my score is lower.} \\
\\ 
\texttt{The maximum combined total in these scenarios is 10. To achieve this, the user needs to take either one hook and the bowl, or one hook and one plate, allowing the other player to take the remaining items. However, since the user's own value is higher for hooks, taking one hook and the bowl gives them 5 points, which is better than 4. The other player would then take the remaining hook and plates, which might be better for them if they value hooks and plates.} \\
\\ 
\texttt{But wait, there} \\
\\ 
\\ 
\texttt{The time is up. It is now your turn to submit a secret proposal.} \\
            }
        }
    }
    & \\ \\

    \theutterance \stepcounter{utterance}  
    & \multicolumn{4}{p{0.6\linewidth}}{
        \cellcolor[rgb]{0.9,0.9,0.9}{
            \makecell[{{p{\linewidth}}}]{
                \texttt{\tiny{[P1$\rangle$GM]}}
                \texttt{<think>} \\
\texttt{Okay, let's try to figure out the best way to divide the items between me and the other player. The items are 2 hooks, 2 plates, and 1 bowl. My values are hook:3, plate:1, bowl:2. The other player's values are secret, so I need to make an allocation that maximizes the total score for both of us.} \\
\\ 
\texttt{First, I need to consider possible allocations where our proposals don't overlap. Since I don't know their values, I have to make some assumptions. Let's think about different scenarios.} \\
\\ 
\texttt{If I take both hooks (3+3=6) and the bowl (2), that's 8 points for me. The remaining items are 2 plates. If the other player values plates more, they could take both, but since plates are only worth 1 each to me, maybe they value them higher. However, if they need the bowl, they might not get it if I take it. This could lead to a conflict.} \\
\\ 
\texttt{Alternatively, if I take one hook (3) and the bowl (2), that's 5 points. The other player can take the remaining hook (3) and the two plates (1+1=2). If they value hooks and plates, their total would be 5. The combined total is 10. That seems good.} \\
\\ 
\texttt{Another option: I take one hook (3) and one plate (1), totaling 4. The other player takes the other hook (3), the bowl (2), and the remaining plate (1). Their total would be 6. Combined total is 10 again. But my score is lower here.} \\
\\ 
\texttt{If I take the bowl (2) and both plates (1+1=2), that's 4 points. The other player takes both hooks (3+3=6). Total combined is 10. But again, my score is lower.} \\
\\ 
\texttt{The maximum combined total in these scenarios is 10. To achieve this, the user needs to take either one hook and the bowl, or one hook and one plate, allowing the other player to take the remaining items. However, since the user's own value is higher for hooks, taking one hook and the bowl gives them 5 points, which is better than 4. The other player would then take the remaining hook and plates, which might be better for them if they value hooks and plates.} \\
\\ 
\texttt{But wait, there} \\
            }
        }
    }
    & & \\ \\

    \theutterance \stepcounter{utterance}  
    & & & \multicolumn{2}{p{0.3\linewidth}}{
        \cellcolor[rgb]{0.9,0.9,0.9}{
            \makecell[{{p{\linewidth}}}]{
                \texttt{\tiny{[GM$|$GM]}}
                \texttt{continue} \\
            }
        }
    }
    & & \\ \\

    \theutterance \stepcounter{utterance}  
    & & & \multicolumn{2}{p{0.3\linewidth}}{
        \cellcolor[rgb]{0.9,0.9,0.9}{
            \makecell[{{p{\linewidth}}}]{
                \texttt{\tiny{[GM$|$GM]}}
                \texttt{player was instructed to make a proposal} \\
            }
        }
    }
    & & \\ \\

\end{supertabular}
}

\end{document}
